% Results' zero-shot paragraph before it was condensed on Sep 29 2026 (Rohan)
\textbf{Zero-shot evaluation on unseen maps.} On unseen maps the turn response survives (U-Net directional scores 0.73 to 0.86 per map against 0.79 to 0.91 in distribution; Table~\ref{tab:unseen}).
The appearance does not: the U-Net's scene PSNR falls from 25.2~dB on the training maps to a median 22.3~dB (20.5 to 24.7 per map) and its LPIPS rises from 0.158 to 0.303 (0.229 to 0.401).
Pooled over the 13 unseen maps, one predicted frame in five has an LPIPS above 0.4, against one in 73 on the training maps; such frames concentrate in maps 16, 17, 13 and 1 (32 to 59 percent of their windows), while seven maps stay below 10 percent.
The decoder is not the limit: it reconstructs the unseen maps nearly as well as the training maps (median 27.3 against 28.6~dB, LPIPS 0.05 to 0.09).
Architecture does not change what survives the shift: in distribution the U-Net's and PixArt-$\alpha$'s scene PSNR differs by 0.02~dB and their LPIPS by 0.001, and the three backbones lose the same things (Table~\ref{tab:unseen}; \appref{app:distance}): scene PSNR drops 2.9, 2.8 and 3.3~dB and LPIPS rises by 0.145, 0.126 and 0.136, so across our three backbones the shift looks like a property of the task, not of the architecture or the scale.
It is a step, not a slope: every unseen map has a worse LPIPS than every training map for all three backbones (\appref{app:distance}).
What the model keeps is how actions act on the world: zero-shot on map 7 the weapon fires and the view advances as in the ground truth (Figure~\ref{fig:teaser}); the loss sits in the scene, whose rows fall 2.9~dB while the map-independent HUD rows fall 0.2~dB (29.7 to 29.5).
